\section{Proposta de Solução e Implementação}
\label{sec:proposta_solucao}

A solução é organizada em componentes de representação, avaliação, vizinhança e busca. Os componentes possuem implementações iniciais em Python, usadas para verificar seu comportamento de forma isolada. \pendente{concluir o registro desses componentes no pyoptframe e atualizar esta seção com a arquitetura integrada.}

\subsection{Representação e organização}
\label{subsec:arquitetura}

A instância é armazenada em JSON. Cada turma contém, entre outros campos, o professor atual, os encontros semanais, as salas, a classificação curricular e, quando disponíveis, os domínios de professores e horários. Durante a avaliação, esses dados são convertidos em tabelas de turmas e encontros.

O código está dividido principalmente em:
\begin{itemize}
    \item \texttt{src/eval/}: leitura da instância, inferência de recursos, tratamento de salas, preferências e avaliador de referência;
    \item \texttt{src/solve/}: avaliador leve e implementações manuais de SA para salas, salas e horários, e professores.
\end{itemize}

Horários marcados como fixos não pertencem à lista de turmas flexíveis do solver de horários. Os domínios usados na implementação inicial foram derivados do histórico de 2025. \pendente{substituir os domínios inferidos pelos dados definitivos de setores e habilitações.}

\subsection{Avaliadores}
\label{subsec:avaliador}

O avaliador de referência, implementado em \texttt{src/eval/evaluator.py}, usa pandas e produz uma decomposição por critério. Ele contabiliza atualmente:
\begin{itemize}
    \item conflitos de sala, professor e grupo curricular;
    \item capacidade insuficiente e incompatibilidade de laboratório;
    \item descanso insuficiente e professores abaixo do mínimo H12;
    \item dias trabalhados, janelas, desperdício de capacidade, rodízio e preferência ponderada.
\end{itemize}

As condições estruturais H1--H5 são tratadas principalmente pela construção da instância e pelos domínios dos movimentos. \pendente{adicionar contadores independentes para as condições estruturais ainda não cobertas, incluindo encontros sem sala em H3.}

Para reduzir o custo causado pelas operações de pandas durante a busca, \path{src/solve/direct_objective.py} implementa a mesma avaliação diretamente sobre o JSON. Esse avaliador é mais leve, mas ainda recalcula a solução completa a cada movimento; não se trata de uma avaliação delta. Testes automatizados verificam a equivalência entre os dois avaliadores para os casos cobertos.

O score usado na configuração inicial é:
\begin{equation}
    F(s)=10^6\,V_{\mathrm{hard}}(s)+\sum_j\Phi_j(s),
\end{equation}
onde $V_{\mathrm{hard}}$ é a soma dos sete contadores fortes implementados e os critérios suaves disponíveis entram com peso unitário. Essa escala serve para testar a prioridade das violações fortes, mas não corresponde aos pesos finais da modelagem.

\subsection{Movimentos implementados}
\label{subsec:vizinhancas}

Foram implementados três tipos de alteração:
\begin{enumerate}
    \item \textbf{sala}: altera a sala de um encontro para outra candidata compatível com a capacidade disponível e a classificação de laboratório;
    \item \textbf{horário}: substitui os encontros de uma turma flexível por outro padrão pertencente ao domínio histórico da turma;
    \item \textbf{professor}: realoca uma turma para outro nome presente no domínio histórico de professores observado para seu setor.
\end{enumerate}

\pendente{implementar a troca simultânea de professores, a troca de salas entre turmas e a composição dos três tipos de decisão em uma mesma busca.}

\subsection{Implementação inicial de Simulated Annealing}
\label{subsec:sa_solver}

Os solvers atuais implementam o SA diretamente em Python. Eles partem da alocação histórica contida na instância, usam um número fixo de iterações e reduzem linearmente a temperatura:
\begin{equation}
    T(i)=\max\left\{1,T_0\left(1-\frac{i}{N}\right)\right\},
\end{equation}
onde $N$ é o total de iterações. O solver conjunto de salas e horários escolhe um movimento de horário com probabilidade $0{,}6$ e um movimento de sala com probabilidade $0{,}4$. O solver de professores executa apenas realocações individuais.

O fluxo comum é apresentado no Algoritmo~\ref{alg:sa_atual}.

\begin{algorithm}[H]
\SetAlgoLined
\KwIn{Instância $I$, temperatura $T_0$, iterações $N$ e semente}
\KwOut{Melhor solução encontrada segundo $F$}
$s\gets\text{cópia da alocação inicial de }I$\;
$s^*\gets s$\;
\For{$i\gets0$ \textbf{até} $N-1$}{
    sortear um movimento permitido e aplicá-lo provisoriamente\;
    $\Delta\gets F(s')-F(s)$\;
    $T\gets\max\{1,T_0(1-i/N)\}$\;
    \eIf{$\Delta\leq0$ \textbf{ou} $\operatorname{Random}(0,1)<\exp(-\Delta/T)$}{
        $s\gets s'$ e atualizar $s^*$ se necessário\;
    }{
        desfazer o movimento\;
    }
}
\Return{$s^*$}\;
\caption{Estrutura da implementação inicial de SA}
\label{alg:sa_atual}
\end{algorithm}

O algoritmo retorna a melhor solução segundo o score, que pode continuar inviável quando nenhuma solução viável é alcançada durante a execução.

\subsection{Integração com pyoptframe}

A implementação no \textit{pyoptframe} utiliza um construtivo, um avaliador e as estruturas de vizinhança correspondentes como componentes registrados no motor do framework. \pendente{implementar a integração, indicar as classes e identificadores empregados e descrever a configuração de SA, ILS e/ou VNS efetivamente comparada.}
